% ============================================================
%  ANSWER KEY — Practice 7: Partitions & Inclusion-Exclusion
% ============================================================

\section*{Վարժություն 7. Տրոհումներ և Ներառման-բացառման սկզբունքը}

\subsubsection*{Խնդիր 1}
Տրոհել 5-տարրանոց բազմությունը 3 ոչ դատարկ դասերի. Ստիռլինգի թիվ $S(5,3)$:

$S(5,3) = S(4,2)+3\cdot S(4,3) = 7 + 3\cdot 6 = \boxed{25}$:

\subsubsection*{Խնդիր 2}
Բաշխել $m$ տարբեր գնդակները $n$ տարբեր արկղերի մեջ այնպես, որ ոչ մի արկղ դատարկ չմնա ($m\to n$ սյուրյեկցիաներ):

Ըստ ներառման-բացառման սկզբունքի՝
\[
  \boxed{n!\,S(m,n) = \sum_{k=0}^{n}(-1)^k\binom{n}{k}(n-k)^m}
\]

\subsubsection*{Խնդիր 3}
Թվեր $\{1,\dots,1000\}$ բազմությունից, որոնք չեն բաժանվում 2-ի, 3-ի կամ 5-ի:

Դիցուք $A_d = $ $d$-ի բազմապատիկները: Ըստ ներառման-բացառման սկզբունքի՝
\begin{align*}
|A_2|&=500,\; |A_3|=333,\; |A_5|=200,\\
|A_6|&=166,\; |A_{10}|=100,\; |A_{15}|=66,\; |A_{30}|=33:
\end{align*}
$|A_2\cup A_3\cup A_5|=500+333+200-166-100-66+33=734$:

$\boxed{1000-734=266}$

\subsubsection*{Խնդիր 4}
4-նիշ թվեր, որոնք պարունակում են գոնե մեկ 7 թվանշան:

Բոլոր 4-նիշ թվերը՝ 9000: Այն թվերը, որոնք չունեն 7. առաջին թվանշանը՝ $8$ ընտրություն, մնացած յուրաքանչյուրը՝ $9$ ընտրություն $= 8\times 9^3 = 5832$:

$\boxed{9000-5832=3168}$

\subsubsection*{Խնդիր 5}
Ամբողջ թվեր $\{1,\dots,16500\}$ բազմությունից, որոնք չեն բաժանվում 5-ի:

5-ի բաժանվողներ. $\lfloor 16500/5\rfloor = 3300$:

$\boxed{16500-3300=13200}$

\subsubsection*{Խնդիր 6}
Ամբողջ թվեր $\{1,\dots,16500\}$ բազմությունից, որոնք չեն բաժանվում ո՛չ 5-ի, ո՛չ էլ 3-ի:

$|A_3|=5500$, $|A_5|=3300$, $|A_{15}|=1100$:
$|A_3\cup A_5|=5500+3300-1100=7700$:

$\boxed{16500-7700=8800}$

\subsubsection*{Խնդիր 7}
Ամբողջ թվեր $\{1,\dots,16500\}$ բազմությունից. բաժանվում են 11-ի, բայց չեն բաժանվում 5-ի և չեն բաժանվում 3-ի:

Ունիվերս $S$. 11-ի բազմապատիկներ, $|S|=1500$: Դիցուք $A=$ նաև բաժ. 3-ի (բաժ. 33-ի), $B=$ նաև բաժ. 5-ի (բաժ. 55-ի):
$|A|=500$, $|B|=300$, $|A\cap B|=\lfloor 16500/165\rfloor=100$:
$|A\cup B|=700$:

$\boxed{1500-700=800}$

\subsubsection*{Խնդիր 8}
Ամբողջ թվեր $\{1,\dots,16500\}$ բազմությունից, որոնք չեն բաժանվում 5-ի, 3-ի կամ 11-ի:

$|A_3|=5500,\;|A_5|=3300,\;|A_{11}|=1500$:
$|A_{15}|=1100,\;|A_{33}|=500,\;|A_{55}|=300$:
$|A_{165}|=100$:

$|A_3\cup A_5\cup A_{11}|=5500+3300+1500-1100-500-300+100=8500$:

$\boxed{16500-8500=8000}$

\subsubsection*{Խնդիր 9}
30 ուսանող; 20-ը գիտեն անգլերեն, 12-ը՝ ֆրանսերեն, 6-ը՝ երկուսն էլ:

$|E\cup F|=20+12-6=26$: Չգիտեն ոչ մեկը.

$\boxed{30-26=4}$

\subsubsection*{Խնդիր 10}
Եռանկյուններ $n=8$ կողմով եռանկյունաձև ցանցում:

Դեպի վեր ուղղված ($s$ կողմով, $1\le s\le 8$). $\sum_{s=1}^{8}\binom{10-s}{2} = 36+28+21+15+10+6+3+1=120$:

Դեպի վար ուղղված ($s$ կողմով, $1\le s\le 4$). $\sum_{s=1}^{4}\binom{10-2s}{2}=28+15+6+1=50$:

$\boxed{120+50=170}$

\subsubsection*{Խնդիր 11}
Երեք գորգ՝ յուրաքանչյուրը 3~մ$^2$ մակերեսով, 6~մ$^2$ մակերեսով սենյակում: Ապացուցել, որ երկուսը հատվում են $\ge \frac{1}{2}$~մ$^2$ մակերեսով:

\textit{Ապացույցի ուրվագիծ.}
Գորգերի ընդհանուր մակերեսը $= 9$~մ$^2$, սենյակի մակերեսը $= 6$~մ$^2$, ուստի ընդհանուր հատումը (հաշված բազմապատիկությամբ) $\ge 9-6=3$~մ$^2$ է: Կան $\binom{3}{2}=3$ զույգ գորգեր: Ըստ Դիրիխլեի սկզբունքի՝ առնվազն մեկ զույգը հատվում է $\ge 3/3 = 1$~մ$^2 \ge \frac{1}{2}$~մ$^2$ մակերեսով: \qed

\subsubsection*{Խնդիր 12}
$x_1+x_2+x_3=40$ հավասարման ամբողջ լուծումները $4\le x_1\le 15$, $9\le x_2\le 18$, $5\le x_3\le 16$ պայմաններով:

Տեղադրենք $y_i = x_i - \text{ստորին սահման}$. $y_1+y_2+y_3=22$, որտեղ $0\le y_1\le 11$, $0\le y_2\le 9$, $0\le y_3\le 11$:

Առանց սահմանափակումների. $\binom{24}{2}=276$:
Խախտումները ըստ ներառման-բացառման սկզբունքի. $|A_1|+|A_2|+|A_3|-(|A_1\cap A_2|+|A_1\cap A_3|+|A_2\cap A_3|)+|A_1\cap A_2\cap A_3| = (66+91+66)-(1+0+1)+0 = 221$:

$\boxed{276-221=55}$